\section{Commutators of long entries}\label{sec:long-entries}

The phase property will let us distribute a long word among several
coordinate blocks. To do this inside the Steinberg group, we must control
the extra factors produced by moving its roots. The main point of this
section is that these factors lie in an abelian subgroup on which the
matrix map is injective. Opposite roots require particular care: their
commutation cannot be inferred just from their matrix images.

Fix an odd integer $L$ and a threshold $N\ge1$ from
Lemma~\ref{lem:phases}. Let $k$ bound the lengths of the defining word
rules, and fix
\begin{equation}\label{eq:long-threshold}
 M\ge 2N+k+1.
\end{equation}
All groups in this section are stable groups over $R$. We identify a
nonzero word class with its basis element in $R$, and write $|w|$ for its
length. Write $E_{ij}$ for the matrix unit in position $(i,j)$.
Phases always take values in $\Z/L$.

\subsection{Opposite roots and movement of a cut}

\begin{lemma}\label{lem:opposite}
Suppose $a,c\in R$ satisfy $ca=0$. For distinct indices $i,h,j$ and
every $b\in R$,
\begin{equation}\label{eq:cut-movement}
 [x_{ij}(ab),x_{ji}(c)]
   =[x_{ih}(a),x_{hi}(bc)].
\end{equation}
In particular, $[x_{ij}(a),x_{ji}(c)]$ is independent of $j\ne i$.
It commutes with both defining roots, has square $1$, and has matrix
image $\Id+acE_{ii}$.
\end{lemma}

\begin{proof}
Apply Lemma~\ref{lem:commutator} to
\[
 X=x_{ih}(a),\qquad Y=x_{hj}(b),\qquad Z=x_{ji}(c).
\]
Then $A=[X,Y]=x_{ij}(ab)$ and $B'=[Y,Z]=x_{hi}(bc)$.
The Steinberg relations give $[A,X]=[B',Z]=1$.
They also give $[X,Z]=1$ because $ca=0$, and $[A,B']=1$
because $(bc)(ab)=b(ca)b=0$. This proves
\eqref{eq:cut-movement}.

Taking $b=1$ proves independence of the auxiliary index. Express the
commutator using an index $h$ different from $i,j$. Each of its two
factors then commutes with $x_{ij}(a)$ and with $x_{ji}(c)$: the only
potential nonzero consecutive-root coefficient is $ca$. Thus the
commutator commutes with both defining roots. Since each root has
square $1$, expanding $[x_{ij}(a)^2,x_{ji}(c)]$ shows
that its square is $1$. Finally, multiplication of the four elementary
matrices, using $ca=0$, gives $\Id+acE_{ii}$.
\end{proof}

\begin{lemma}\label{lem:long-commutator}
For each coordinate $i$ and every nonzero word class $w$ satisfying
$|w|\ge M$ and $\sigma(w)\ne\tau(w)$, there is an element
$H_i(w)\in\St(R)$ with the following property. If $p,q$ are nonzero
words with
\[
 |p|,|q|\ge N,\qquad |p|+|q|\ge M,
 \qquad \sigma(p)\ne\tau(q),
\]
then, for every $j\ne i$,
\begin{equation}\label{eq:long-opposite}
 [x_{ij}(p),x_{ji}(q)]
 =\begin{cases}
    H_i(w),& pq\text{ represents the nonzero class }w,\\
    1,& pq=0\text{ in }R.
   \end{cases}
\end{equation}
The element $H_i(w)$ has square $1$ and matrix image
$\Id+wE_{ii}$.
\end{lemma}

\begin{proof}
The phase mismatch gives $qp=0$, so
Lemma~\ref{lem:opposite} applies. We first show that the commutator
does not change when the cut between $p$ and $q$ moves, provided both
new parts still have length at least $N$.

For a move to the right, write the original word as $abc$, with
$a=p$ and $bc=q$, and with $|c|\ge N$. The word $c$ is nonzero and
has $\tau(c)=\tau(q)$. Hence $ca=0$, and
\eqref{eq:cut-movement}, together with auxiliary-index independence,
identifies the commutators for the cuts $a\mid bc$ and $ab\mid c$.
This equality remains valid if $ab=0$, in which case the commutator
is $1$. Otherwise $\sigma(ab)=\sigma(a)$, so the new cut retains
the mismatch. For a move to the left, write $p=ab$ and $q=c$ with
$|a|\ge N$. Now $a$ is nonzero with $\sigma(a)=\sigma(p)$, so
$ca=0$ again, and the same identity compares $ab\mid c$ with
$a\mid bc$. If the new suffix is zero we again obtain $1$;
otherwise its end phase is unchanged. This proves invariance under
every allowed movement of the cut. If a part becomes zero, the
equalities already obtained show that the original commutator was
$1$, and the argument can stop.

Next consider a single defining replacement inside a word of length
$n\ge M$. The allowed cut positions are the integers in $[N,n-N]$.
A rule of length at most $k$ can straddle at most $k-1$ of these
positions, whereas \eqref{eq:long-threshold} supplies more than $k$
allowed positions. We can therefore move the cut so that the entire
replacement lies in one part. That part has the same value in $R$
before and after the replacement. If it is nonzero, its phases are
unchanged because phases depend only on the word class. The other
part is unchanged. The commutator consequently survives the
replacement with the same value and mismatch, unless a zero part has
already proved the original commutator to be $1$.

If the concatenation is nonzero, every part of every equivalent word
is nonzero. Cut movement and the preceding replacement argument show
that the commutator depends only on the class of the concatenation.
Conversely, a nonzero class $w$ as in the statement can be cut into
two parts of lengths at least $N$. The subword assertion of
Lemma~\ref{lem:phases} makes their start and end phases
$\sigma(w)$ and $\tau(w)$, respectively. Define $H_i(w)$ by any
such cut. The independence just proved makes this definition
unambiguous.

If the concatenation represents zero, the criterion
\eqref{eq:word-zero-criterion} gives a sequence of replacements
ending at a word that contains a zero-declared subword. Follow this
sequence with the commutator. Either a part has already become zero,
or the cut can be moved once more to put the final zero-declared
subword wholly in one part. That part is zero in $R$, and the
commutator is $1$. This proves the zero case as well. The order and
matrix-image assertions now follow from
Lemma~\ref{lem:opposite}.
\end{proof}

\subsection{A subgroup detected by its matrix entries}

Let $S$ be a finite set of coordinates, with a phase
$P(i)\in\Z/L$ assigned to each $i\in S$. An entry $w$ from $i$ to
$j$ will be called \emph{consistent} if it is a nonzero word class
of length at least $M$ and
\[
 \sigma(w)=P(i),\qquad \tau(w)=P(j).
\]
We call it a \emph{wrong-row entry} if instead
\begin{equation}\label{eq:wrong-row}
 |w|\ge M,\qquad \sigma(w)\ne P(i),\qquad \tau(w)=P(j).
\end{equation}
For $i\ne j$, a wrong-row entry defines the root $x_{ij}(w)$.
For $i=j$, it defines $H_i(w)$ by
Lemma~\ref{lem:long-commutator}. Let $U(S,P)$ be the subgroup
generated by all these elements.

\begin{lemma}\label{lem:wrong-row}
The subgroup $U(S,P)$ is abelian of exponent at most $2$, and the
matrix map is injective on it. Every root with a consistent entry
normalizes $U(S,P)$.
\end{lemma}

\begin{proof}
Write $U=U(S,P)$. First consider an off-diagonal generator
$x_{ij}(a)$ of $U$ and an off-diagonal root $x_{ef}(b)$ whose entry
is either wrong-row or consistent. If $j=e$ and $i\ne f$, their
commutator is $x_{if}(ab)$. The product is zero when $b$ is
wrong-row, since $\tau(a)=P(j)\ne\sigma(b)$. When $b$ is
consistent and $ab\ne0$, the product inherits $\sigma(a)$ and
$\tau(b)$, so it is a wrong-row entry from $i$ to $f$.
If $i=f$ and $j\ne e$, the reverse product $ba$ is zero, because
$\tau(b)=P(i)\ne\sigma(a)$, and the roots commute.

If the two positions are opposite, so that $e=j$ and $f=i$, the
same mismatch $\sigma(a)\ne\tau(b)$ allows us to use
\eqref{eq:long-opposite}. The commutator is $H_i(ab)$ when
$ab\ne0$, and is $1$ when $ab=0$. In the first case $ab$ is a
wrong-row entry from $i$ to itself. In particular the commutator is
$1$ whenever $b$ is wrong-row. All remaining pairs of positions
commute by the defining Steinberg relations. This proves both the
required commutations between off-diagonal generators of $U$ and
their required behavior under consistent roots.

Now take a diagonal generator $H_i(w)$. Represent it as a
commutator through an auxiliary coordinate $h\notin S$. To compare
it with a root $x_{ef}(b)$ on $S$, choose $h$ outside its indices
as well. If neither index is $i$, the supports are disjoint and the
elements commute. If one index is $i$, the root lies in a strip
between $\{i,h\}$ and its complement. By
Lemma~\ref{lem:strip}, conjugation by $H_i(w)$ acts on that strip
by its matrix $\Id+wE_{ii}$. This matrix is its own inverse,
since $w^2=0$ by the phase mismatch. Consequently, for $f\ne i$
and $e\ne i$,
\begin{align*}
 H_i(w)x_{if}(b)H_i(w)^{-1}&=x_{if}(b+wb),\\
 H_i(w)x_{ei}(b)H_i(w)^{-1}&=x_{ei}(b+bw).
\end{align*}
The incoming product $bw$ is always zero, because
$\tau(b)=P(i)\ne\sigma(w)$. The outgoing product $wb$ is zero
when $b$ is wrong-row; when $b$ is consistent and $wb\ne0$, it is
a wrong-row entry from $i$ to $f$. These formulas prove that
diagonal generators commute with off-diagonal generators of $U$,
and that their commutators with consistent roots belong to $U$.

Two diagonal generators at different coordinates can be represented
using disjoint pairs of coordinates, and hence commute. For two at
the same coordinate $i$, write the second as
\[
 H_i(w')=[x_{it}(a),x_{ti}(c)],\qquad ac=w',\quad |a|,|c|\ge N,
\]
using $t$ distinct from the auxiliary coordinate for the first.
Here $\sigma(a)=\sigma(w')\ne P(i)$ and
$\tau(c)=\tau(w')=P(i)$. The phase rule therefore gives
$wa=0$ and $cw=0$. Applying the same strip action to the two
displayed factors shows that $H_i(w)$ fixes both, and thus commutes
with $H_i(w')$. The use of parts of length at least $N$ is sufficient
for these two applications of the phase rule.

We have proved that all generators of $U$ commute. They have square
$1$ by root addition and Lemma~\ref{lem:long-commutator}. Their
matrix images are $\Id+wE_{ij}$, indexed by positions and nonzero
word classes satisfying \eqref{eq:wrong-row}. Any two added matrices
multiply to zero: a potentially nonzero product requires matching
positions, and then the end phase of the first entry is the designated
row phase of the second, which differs from its start phase. Thus a
product of distinct generators has matrix image
\[
 \Id+\sum_{(i,j,w)}wE_{ij}.
\]
The word classes form a basis of $R$, so this is the identity only
when the sum is empty. This proves injectivity on $U$.
Finally, the commutator calculations show that conjugation by any
consistent root sends every generator of $U$ into $U$. Such a root
is its own inverse, so it normalizes $U$.
\end{proof}
